\chapter{The FX Market}\label{ch:m2:the-fx-market}

In April 2025 the world's banks and dealers traded, on an average day, USD~9.6
trillion of foreign exchange, 28\% more than three years earlier. There is no
exchange where this happens, no single price and no closing time. The
market is a web of dealers who quote to each other and to their clients, a few
electronic order books that most of the world watches, dozens of platforms
that compete to connect them, and a settlement system that makes sure each side
of a trade delivers its currency. The US dollar is on one side of almost nine
trades in ten. This chapter describes the market's conventions, its
structure, the venues in which it trades and what the central banks' survey of
it says; the following chapters take its instruments one by one.

\section{Currencies, pairs and conventions}

\begin{definition}[Currency pair, base and quote currency, pip]\label{def:m2:the-fx-market:pair}
A \emph{currency pair}\index{currency pair} is written as two three-letter codes,
BASEQUOTE, and its rate $S$ is the number of units of the second, the
\emph{quote currency}\index{quote currency}, per unit of the first, the
\emph{base currency}\index{base currency}. A \emph{pip}\index{pip} is the
conventional unit of a rate's movement: 0.0001 for most pairs, 0.01 for pairs
quoted in yen.
\end{definition}

EURUSD at 1.1464 means 1.1464 dollars per euro; USDJPY at 156.87 means 156.87
yen per dollar. Which currency is the base is a convention, not a choice: the
Federal Reserve's daily rates are quoted in units of currency per dollar except
for the euro, sterling, and the Australian and New Zealand dollars, which are
quoted in dollars per unit; the European Central Bank quotes every rate per euro,
so EURGBP is pounds per euro. A trader who \emph{buys} EURUSD buys euros and
sells dollars; a quote of 1.1462/1.1464 means that the dealer buys euros at the
bid, 1.1462, and sells them at the ask, 1.1464, a spread of two \omterm{def:m2:the-fx-market:pair}{pips}.

\begin{definition}[Spot value date, cross rate]\label{def:m2:the-fx-market:spot}
The \emph{spot value date}\index{spot value date} of a trade is the day on which
the two currencies are exchanged: by convention the second business day after
the trade, counting only days that are business days in both currencies, and
the first for dollar--Canadian dollar. A \emph{cross rate}\index{cross rate} is
the rate between two currencies neither of which is the dollar, such as EURJPY,
traded directly or through the two dollar pairs.
\end{definition}

The rule has a detail every system must encode. Federal Reserve economists
described it in 2003 with dollar--yen: a trade on a Tuesday settles on
Thursday; if Wednesday is a holiday in Japan, on Friday; but if Wednesday is a
holiday in the United States and not in Japan, still on Thursday, because a
dollar holiday on the intermediate day does not delay a dollar pair
(\cref{fig:m2:the-fx-market:spotdate}). A day ends at 17:00 in New York: a
trade just after it is dated the next day and settles a day later. Everything
beyond spot is a forward, the subject of \cref{ch:m2:fx-swaps-forwards-and-the-cross-currency-basis}.

\begin{omfigure}
\centering
\begin{tikzpicture}[font=\small,x=1.9cm,y=0.75cm]
\foreach \d/\n in {0/Tue,1/Wed,2/Thu,3/Fri} {\node[font=\footnotesize\bfseries] at (\d,0.9) {\n};}
\node[anchor=east,font=\footnotesize] at (-0.45,0) {no holiday};
\node[anchor=east,font=\footnotesize] at (-0.45,-1) {Japan holiday Wed};
\node[anchor=east,font=\footnotesize] at (-0.45,-2) {US holiday Wed};
\node[anchor=east,font=\footnotesize] at (-0.45,-3) {US holiday Thu};
\foreach \r/\v/\h in {0/2/-1,-1/3/1,-2/2/1,-3/3/2} {
  \foreach \d in {0,1,2,3} {\draw[gray!50] (\d-0.4,\r-0.35) rectangle (\d+0.4,\r+0.35);}
  \node[omDef,font=\footnotesize] at (0,\r) {trade};
  \fill[omThm!25] (\v-0.4,\r-0.35) rectangle (\v+0.4,\r+0.35);
  \node[omThm,font=\footnotesize\bfseries] at (\v,\r) {value};
}
\node[omProp,font=\footnotesize] at (1,-1) {JP closed};
\node[omProp,font=\footnotesize] at (1,-2) {US closed};
\node[omProp,font=\footnotesize] at (2,-3) {US closed};
\end{tikzpicture}

\omcaption{\omterm{def:m2:the-fx-market:spot}{Spot value dates} of a dollar--yen trade done on a Tuesday. A Japanese
holiday on the intermediate day pushes settlement to Friday; a US holiday on it
does not; a US holiday on the value date itself does. The rule as described by
Federal Reserve economists in 2003; holidays
hypothetical.}\label{fig:m2:the-fx-market:spotdate}
\end{omfigure}

A \omterm{def:m2:the-fx-market:spot}{cross rate} can always be built from the two dollar pairs, and the
construction fixes which side of each leg is used.

\begin{proposition}[Triangulated cross]\label{prop:m2:the-fx-market:cross}
With EURUSD quoted $b_1/a_1$ and USDJPY $b_2/a_2$, the synthetic EURJPY is
$b_1b_2 / a_1a_2$. With EURUSD $b_1/a_1$ and GBPUSD $b_3/a_3$, the synthetic
EURGBP is $(b_1/a_3)/(a_1/b_3)$. In each case the relative spread of the cross
is close to the sum of the legs' relative spreads.
\end{proposition}

\begin{proof}
To sell euros for yen through the dollar, sell euros at the EURUSD bid $b_1$ and
sell those dollars for yen at the USDJPY bid $b_2$: $b_1b_2$ yen per euro. To buy
euros, pay the two asks. For EURGBP, selling euros gives $b_1$ dollars, and
buying pounds with dollars costs $a_3$ dollars per pound. Relative spreads add to
first order because $\log(a_1a_2) - \log(b_1b_2) = (\log a_1 - \log b_1) +
(\log a_2 - \log b_2)$.
\end{proof}

\begin{example}[Three crosses]\label{ex:m2:the-fx-market:crosses}
With EURUSD 1.1462/1.1464, USDJPY 156.86/156.88 and GBPUSD 1.3371/1.3373, the
synthetic EURJPY is 179.7929/179.8472, 5.43 \omterm{def:m2:the-fx-market:pair}{pips} wide, a relative spread of
3.02 basis points against 1.74 and 1.27 for the legs; EURGBP is
0.85710/0.85738, 2.78 \omterm{def:m2:the-fx-market:pair}{pips}; GBPJPY 209.7375/209.7956, 5.81 \omterm{def:m2:the-fx-market:pair}{pips}. A dealer that
quotes a cross directly can quote it tighter than the synthetic, but not
crossed with it: a direct EURJPY bid above 179.8472 could be sold to and bought
back through the legs at a profit.
\end{example}

\section{Dealers, tiers and the over-the-counter structure}

The market was built by banks trading with each other and with their clients,
and still has that shape. The 2025 survey counts 46\% of turnover as trades
between reporting dealers and 50\% as trades with other financial institutions,
such as smaller banks, asset managers, hedge funds and principal trading firms;
non-financial companies are the rest. Four places, the United Kingdom, the
United States, Singapore and Hong Kong, book three quarters of it, London alone
38\%.

A large dealer quotes thousands of clients in hundreds of pairs, and nets their
flows against each other before it goes to the market: a client selling euros
and another buying them within the same second leave the dealer with no risk,
and neither trade reaches an interdealer venue. This internalisation (One Quant
Book 1, chapter 10) is why the dealer market is concentrated: the more flow a dealer sees,
the more it can net, the tighter it can quote and the more flow it attracts.
Smaller banks and regional dealers, the second tier, take prices from the large
ones and quote their own clients on top. Principal trading firms, which trade
their own capital with fast algorithms, have become market makers too,
streaming prices to banks and to platforms (\cref{ch:m2:fx-spot-microstructure}).

\begin{omfigure}
\centering
\begin{tikzpicture}[font=\small,
  box/.style={draw,thick,rounded corners=2pt,align=center,minimum height=0.9cm,minimum width=2.6cm,inner sep=3pt},
  arr/.style={{Stealth[length=2mm]}-{Stealth[length=2mm]},thick}]
\node[box,draw=omThm,fill=omThm!10] (pv) at (5,3.2) {primary venues\\(central limit order books)};
\node[box,draw=omDef,fill=omDef!8] (d1) at (1.6,1.4) {large dealers};
\node[box,draw=omDef,fill=omDef!8] (ptf) at (8.4,1.4) {principal\\trading firms};
\node[box,draw=gray,fill=gray!8] (ecn) at (5,0.2) {ECNs and\\multi-dealer platforms};
\node[box,draw=omProp,fill=omProp!8] (sdp) at (1.6,-1.4) {single-dealer\\platforms};
\node[box,draw=omProp,fill=omProp!8] (cl) at (5,-3.0) {clients: funds, companies,\\smaller banks};
\node[box,draw=omThm,fill=omThm!8] (pb) at (8.4,-1.4) {prime broker\\(credit)};
\draw[arr] (d1) -- (pv); \draw[arr] (ptf) -- (pv);
\draw[arr] (d1) -- (ecn); \draw[arr] (ptf) -- (ecn);
\draw[arr] (d1) -- (sdp); \draw[arr] (sdp) -- (cl);
\draw[arr] (ecn) -- (cl);
\draw[arr,dashed] (cl) -- (pb); \draw[arr,dashed] (pb) -- (ptf);
\end{tikzpicture}

\omcaption{The structure of the spot market. Dealers and principal trading
firms trade with each other on the \omterm{def:m2:the-fx-market:venue}{primary venues} and on other electronic
networks, and with clients on those networks and on their own \omterm{def:m2:the-fx-market:sdp}{single-dealer
platforms}; a prime broker lends a client its name and credit to trade with
dealers it could not otherwise face (dashed). Schematic, after the Bank for
International Settlements' descriptions.}\label{fig:m2:the-fx-market:structure}
\end{omfigure}

\section{Primary venues and electronic networks}

\begin{definition}[Primary venue, electronic communication network]\label{def:m2:the-fx-market:venue}
A \emph{primary venue}\index{primary venue} is one of the interdealer electronic
order books that the market treats as the reference for a pair's price. An
\emph{electronic communication network}\index{electronic communication network}
(ECN) is any other electronic trading platform where several liquidity
providers and takers meet, by order book or by request for quote, anonymously
or with names disclosed.
\end{definition}

For three decades two electronic brokers, Refinitiv Matching, formerly
Reuters, and EBS Market, now owned by CME Group, have been the \omterm{def:m2:the-fx-market:venue}{primary venues} of
the spot market,
central limit order books in which dealers, and since the mid-2000s principal
trading firms, post and hit prices; their prices have been the references for
the rest of the market. The BIS counted more than thirty secondary venues in the
dealer-to-client segment in 2022, and volumes on the \omterm{def:m2:the-fx-market:venue}{primary venues} declining
for more than a decade: dealers internalise more, execution algorithms slice
orders across many venues, and dealers wary of being picked off by faster firms
pushed the venues to add small delays, speed bumps, to incoming orders.
Much price discovery now also happens in currency futures.

\section{Single-dealer platforms and prime brokerage}

\begin{definition}[Single-dealer platform, FX prime brokerage]\label{def:m2:the-fx-market:sdp}
A \emph{single-dealer platform}\index{single-dealer platform} is a dealer's own
electronic system through which its clients see its prices and trade with it.
\emph{FX prime brokerage}\index{FX prime brokerage} is the service by which a
dealer lets a client trade with a set of other dealers in the prime broker's name
and on its credit: each trade is given up to the prime broker, which becomes
the counterparty of the executing dealer and of the client.
\end{definition}

The 2022 survey recorded a shift of seven percentage points of trading towards
direct electronic channels, \omterm{def:m2:the-fx-market:sdp}{single-dealer platforms} and direct price streams,
away from anonymous venues. For a client a \omterm{def:m2:the-fx-market:sdp}{single-dealer platform} offers
personalised prices and, often, the dealer's other services; for a dealer it is
the best way to see its clients' flow and to internalise it.

Prime brokerage is what lets a hedge fund or a principal trading firm act like a
bank. The executing dealer sees only the prime broker's name, so the client
trades anonymously and with many dealers under one credit relationship, one set
of settlements and one report; the prime broker charges a fee and bears the
client's credit risk. The BIS estimated that close to a third of the turnover
with financial customers was prime-brokered in 2019, and \cref{ch:m2:getting-access-fx}
returns to what it takes to get such an arrangement and to the limits that come
with it.

\section{What the turnover surveys say}

Every three years central banks ask the dealers in their jurisdictions to
report their trading in April, and the BIS assembles the results: in 2025, 52
jurisdictions and more than 1\,100 dealers. The survey counts each trade once
across the world, by the location of the dealer's sales desk.

\begin{omfigure}
\begin{tikzpicture}
\begin{axis}[width=11cm,height=4.8cm,ybar,bar width=10pt,
  ylabel={share of turnover (\%)},
  tick label style={font=\small},label style={font=\small},
  xmin=-0.6,xmax=4.6,ymin=0,ymax=60,grid=major,grid style={gray!25},
  xtick={0,1,2,3,4},xticklabels={spot,outright\\forwards,FX\\swaps,options,currency\\swaps},
  xticklabel style={align=center},
  nodes near coords,every node near coord/.append style={font=\scriptsize},
  legend columns=2,legend style={font=\footnotesize,at={(0.5,-0.42)},anchor=north,draw=none,fill=none,/tikz/every even column/.append style={column sep=8pt}},
  table/col sep=comma]
\addplot[fill=gray!40,draw=gray] table[x=k,y=y2022]{figdata/markets-2/14-the-fx-market/instruments.csv};
\addplot[fill=omThm!60,draw=omThm] table[x=k,y=y2025]{figdata/markets-2/14-the-fx-market/instruments.csv};
\legend{April 2022 (USD 7.5 trillion a day),April 2025 (USD 9.6 trillion a day)}
\end{axis}
\end{tikzpicture}

\omcaption{Global FX turnover by instrument, as a share of the daily total, in
April 2022 and April 2025. FX swaps remain the largest instrument, but spot,
forwards and options grew faster. Data: BIS Triennial Central Bank Survey,
2025 release.}\label{fig:m2:the-fx-market:instruments}
\end{omfigure}

Two messages stand out. First, spot is less than a third of the market: most
trading is in FX swaps and forwards, instruments of funding and hedging
(\cref{fig:m2:the-fx-market:instruments}). Second, the dollar is the
market's vehicle (\cref{fig:m2:the-fx-market:currencies}): the ten most
traded pairs all include it, and a Swedish exporter paid in yen often finds it
cheaper to sell yen for dollars and dollars for kronor than to trade the thin
cross.
Shares add up to 200\% because each trade involves two currencies. The survey
also warned that April 2025 was not a quiet month: it followed trade-policy
announcements that moved currencies sharply.

\begin{omfigure}
\begin{tikzpicture}
\begin{axis}[width=11cm,height=4.6cm,ybar,bar width=14pt,
  ylabel={share of trades (\%)},
  tick label style={font=\small},label style={font=\small},
  xmin=-0.6,xmax=6.6,ymin=0,ymax=110,grid=major,grid style={gray!25},ytick={0,20,40,60,80,100},
  xtick={0,1,2,3,4,5,6},xticklabels={USD,EUR,JPY,GBP,CNY,CHF,HKD},
  nodes near coords,every node near coord/.append style={font=\scriptsize},
  table/col sep=comma]
\addplot[fill=omDef!60,draw=omDef] table[x=k,y=y2025]{figdata/markets-2/14-the-fx-market/currencies.csv};
\end{axis}
\end{tikzpicture}

\omcaption{Share of global FX trades involving each currency, April 2025 (the
shares sum to 200\% since every trade has two sides). The Australian and
Canadian dollars, at about 6\% each, are not shown. Data: BIS Triennial Central
Bank Survey, 2025 release.}\label{fig:m2:the-fx-market:currencies}
\end{omfigure}

\begin{dated}{2026-09}{Size and settlement}\label{dat:m2:the-fx-market:survey}
BIS Triennial Survey, April 2025: USD~9.6 trillion a day; spot USD~3 trillion,
outright forwards USD~1.8 trillion, FX swaps USD~4 trillion; dollar on one side
of 89.2\% of trades, euro 28.9\%, yen 16.8\%, sterling 10.2\%, renminbi 8.5\%,
Swiss franc 6.4\%. The survey is repeated every three years. CLS, the settlement
system for the main currencies (\cref{ch:m2:settlement-risk}), settles 18
currencies, on average more than USD~8 trillion a day, for more than 75
settlement members.
\end{dated}

\section{Tutorial: spot dates and crosses}\label{tut:m2:the-fx-market:pairs}

\begin{tutorial}
\textbf{Goal.} Name pairs by convention, compute \omterm{def:m2:the-fx-market:spot}{spot value dates} across two
holiday calendars, and triangulate crosses with their bid and ask.
\textbf{End state:} \cref{fig:m2:the-fx-market:spotdate},
\cref{ex:m2:the-fx-market:crosses} and the numbers of the weekend problem.
\begin{tutsteps}
\item \textbf{The spot date}, with the dollar exception on the intermediate
day.
\omcode{code/firm/fxpairs/firm_fxpairs.py}{35}{50}{Spot value date of a pair across two holiday calendars.}\label{lst:m2:the-fx-market:spot}
\item \textbf{The cross}, each leg on the side the trade forces, and the
arbitrage check of a direct quote.
\omcode{code/firm/fxpairs/firm_fxpairs.py}{53}{81}{A cross rate through the dollar, and the arbitrage test.}\label{lst:m2:the-fx-market:cross}
\item \textbf{Run} \texttt{fx\_demo.crosses()}, \texttt{fx\_demo.spot\_examples()},
\texttt{fx\_demo.the\_cross()} and \texttt{fig\_fx.py}.
\end{tutsteps}
\textbf{What to change next.} Add a holiday calendar from a public source for
the euro area (the TARGET closing days) and Japan, and compute a year of EURJPY
spot dates; then list the days on which the spot date jumps by more than two
business days.
\end{tutorial}

\section{Build: the pair conventions}\label{bld:m2:the-fx-market:fxpairs}

\begin{build}
\bfield{Purpose} Every FX system of the miniature firm starts here: the name
and orientation of a pair, its \omterm{def:m2:the-fx-market:pair}{pip}, its spot date, and the synthetic price of a
cross.
\bfield{Interface} \texttt{pair\_name(a, b, priority)}; \texttt{pip(pair)};
\texttt{spot\_date(trade, pair, holidays)}; \texttt{Quote(bid, ask)};
\texttt{cross\_via\_usd(leg1, leg2, target)}; \texttt{spread\_pips};
\texttt{arbitrage(direct, synthetic)}.
\bfield{Rules} Base-currency priority kept as data; holidays passed in, never
hard-coded; T+2 in both currencies' business days, T+1 for USDCAD, a US
holiday on the intermediate day ignored for dollar pairs; each leg of a cross
on the side the trade forces.
\bfield{Acceptance tests} \texttt{code/firm/fxpairs/tests/}: pair names and \omterm{def:m2:the-fx-market:pair}{pips};
the four published dollar--yen cases and USDCAD; a Friday trade settling on
Tuesday; crosses equal to the products and quotients of the right sides; a
crossed direct quote flagged as arbitrage.
\bfield{Stretch} Real holiday calendars, versioned; the 17:00 New York roll;
value dates for forwards (\cref{ch:m2:fx-swaps-forwards-and-the-cross-currency-basis}); crosses through a
currency other than the dollar.
\end{build}

\begin{omsources}
\item Bank for International Settlements, Triennial Central Bank Survey,
``OTC foreign exchange turnover in April 2025'' and press release, 30 September
2025.
\item Bank for International Settlements, \emph{Quarterly Review}, December 2019
(``Sizing up global foreign exchange markets'') and December 2022 (boxes on
trade execution and the fragmented spot market).
\item A.~Chaboud and J.~Wright, ``Uncovered interest parity: it works, but not
for long'', Federal Reserve Board, International Finance Discussion Paper 752,
2003.
\item Federal Reserve, H.10 foreign exchange rates; FR 3036 survey
instructions; Foreign Exchange Committee, amendments to the 1998 FX
definitions, 2011.
\item CLS Group, company information.
\end{omsources}

\section{Exercises}

\begin{exercise}[$\star$]\label{exo:m2:the-fx-market:1}
EURUSD is 1.1462/1.1464. A client sells EUR~5 million. How many dollars does
it receive, and what is the spread in \omterm{def:m2:the-fx-market:pair}{pips} and in basis points of the mid?
\end{exercise}

\begin{exercise}[$\star$]\label{exo:m2:the-fx-market:2}
Name the pairs formed by the dollar and the yen, the euro and sterling, the
Australian dollar and the yen, and give the \omterm{def:m2:the-fx-market:pair}{pip} of each.
\end{exercise}

\begin{exercise}[$\star$]\label{exo:m2:the-fx-market:3}
A USDJPY trade is done on Tuesday 6 October 2026. Give its spot date with no
holidays, with a Japanese holiday on Wednesday, and with a US holiday on
Thursday. A USDCAD trade on the same day?
\end{exercise}

\begin{exercise}[$\star\star$]\label{exo:m2:the-fx-market:4}
From EURUSD 1.1462/1.1464 and GBPUSD 1.3371/1.3373, derive the synthetic EURGBP
and its spread in \omterm{def:m2:the-fx-market:pair}{pips}.
\end{exercise}

\begin{exercise}[$\star\star$]\label{exo:m2:the-fx-market:5}
Why do the survey's currency shares add up to 200\%, and what does the dollar's
89\% say about how a cross such as SEKJPY is usually traded?
\end{exercise}

\begin{exercise}[$\star\star$]\label{exo:m2:the-fx-market:6}
A hedge fund with a prime broker executes USD~50 million of EURUSD with a
dealer it has no agreement with. Who is the dealer's counterparty, and who
bears the fund's credit risk?
\end{exercise}

\begin{exercise}[$\star\star\star$]\label{exo:m2:the-fx-market:7}
\emph{Coding.} With \texttt{cross\_via\_usd}, give the synthetic GBPJPY and its
relative spread, and check that it is close to the sum of the legs' relative
spreads.
\end{exercise}

\begin{exercise}[$\star\star\star$]\label{exo:m2:the-fx-market:8}
\emph{Find the flaw.} ``FX turnover is USD~9.6 trillion a day, so an exporter
can sell any amount of any currency at the screen price.'' Correct it.
\end{exercise}

\section{Problem: The Cross}

\begin{problem}[{Weekend problem --- the most a client should pay for a cross}]\label{pb:m2:the-fx-market:1}
A client wants to buy EUR~10 million against yen. The market shows EURUSD
1.1462/1.1464 and USDJPY 156.86/156.88. Two dealers quote EURJPY directly:
dealer~A 179.80/179.84, dealer~B 179.85/179.87.

\medskip\noindent\textbf{Part I --- The synthetic cross.}
\begin{enumerate}
\item Which side of each leg does the client hit to buy euros with yen?
\item Give the synthetic EURJPY bid and ask.
\item Give its spread in \omterm{def:m2:the-fx-market:pair}{pips} and in basis points of the mid.
\item Give the legs' relative spreads and compare their sum with the cross's.
\item How many dollars pass through the trade if the client buys synthetically?
\end{enumerate}

\medskip\noindent\textbf{Part II --- The direct quotes.}
\begin{enumerate}[resume]
\item Is dealer~A's quote consistent with the legs?
\item Is dealer~B's? What would an arbitrageur do?
\item Give the arbitrage profit on EUR~10 million, in yen and in dollars.
\item Why might such a quote exist for a moment, and why not for long?
\item Where should the client buy?
\end{enumerate}

\medskip\noindent\textbf{Part III --- Costs beyond the quote.}
\begin{enumerate}[resume]
\item Why does trading the legs cost the client more than the arithmetic?
\item When is the direct cross cheaper than the synthetic?
\item How does settlement differ between the two routes?
\item What changes if Wednesday is a holiday in Japan and the trade is on Tuesday?
\item Why can a dealer quote the cross tighter than the synthetic?
\end{enumerate}

\medskip\noindent\textbf{Part IV --- Judgement.}
\begin{enumerate}[resume]
\item What is the danger in quoting the cross off stale legs?
\item Why do thin crosses trade through the dollar?
\item How would a large order in the cross move the legs?
\item State the \emph{named result}: the widest spread the client should accept
for EURJPY today.
\item In one sentence: what bounds the price of any cross?
\end{enumerate}
\end{problem}

\section{Interview questions}

\begin{interviewq}[$\star$ \iqroles{trader, developer}]\label{iq:m2:the-fx-market:1}
What does EURUSD 1.1462/1.1464 mean, and which side does a client who buys euros
hit?
\end{interviewq}

\begin{interviewq}[$\star$ \iqroles{developer}]\label{iq:m2:the-fx-market:2}
How do you compute the spot date of an FX trade? What edge cases would you test?
\end{interviewq}

\begin{interviewq}[$\star\star$ \iqroles{trader, researcher}]\label{iq:m2:the-fx-market:3}
Derive the bid and ask of EURGBP from EURUSD and GBPUSD.
\end{interviewq}

\begin{interviewq}[$\star\star$ \iqroles{trader, bank}]\label{iq:m2:the-fx-market:4}
Why is the FX market over the counter rather than on an exchange, and how has
electronic trading changed it?
\end{interviewq}

\begin{interviewq}[$\star\star$ \iqroles{bank}]\label{iq:m2:the-fx-market:5}
What does a prime broker do for an FX client, and what risks does it take?
\end{interviewq}

\begin{interviewq}[$\star\star\star$ \iqroles{developer, trader}]\label{iq:m2:the-fx-market:6}
Design the component that publishes synthetic cross prices for a hundred
crosses from the dollar legs, with a latency budget of a few microseconds.
\end{interviewq}
